% generated from baselines/simple_vqa/method.py and baselines/rbench/prompt.py
\paragraph{Simple VQA prompt.}
\begin{lstlisting}[style=promptstyle]
System: You are a concise one-shot video evaluator. Apply the rubric and emit only the function call.

User: Evaluate this embodied-action video against the supplied goal in one direct pass.
Do not generate a plan, subgoal list, dependency graph, or gate-by-gate answers.
Do not infer hidden events. Return one structured function call containing the
two final rubric scores and concise observable evidence.

TASK COMPLETION (0-3)
Judge only whether the requested task was attempted and completed. Ignore
visual glitches, clipping, malformed hands, floating, teleportation, jitter,
and impossible physics when the macro-intent and requested outcome remain
recognizable. Use the whole goal, including multiple requested outcomes and
their necessary ordering, but report one holistic score.
- 0: no recognizable task-directed attempt, or behavior targets the wrong
  source object, tool, or mechanism.
- 1: a correct attempt begins, but the characteristic high-level action is
  not substantially carried through end to end.
- 2: the high-level action is substantially performed, but the required final
  object relation, quantity, or state is materially missed or incomplete.
- 3: the requested final state/outcomes are visibly achieved. Allow slight
  imprecision that does not materially change the stated goal.
Judge final-state evidence more strictly than motion appearance. A prerequisite
or ordering failure should reduce Task when it prevents the requested sequence,
but do not invent customary ordering requirements not implied by the goal.

PHYSICAL PLAUSIBILITY (0-4 or null)
Judge independently of Task success. Return null only when no relevant object
or embodied movement can be inspected; a static/no-attempt clip is Task 0 and
Physics null, not Physics 4. Otherwise assess the visible interaction as a
causal physical event:
- 0: object identity, count, colour/material, or traceable continuity/path
  breaks without a licensed visible transformation.
- 1: an object moves or changes without visible contact, support, or another
  valid preceding trigger.
- 2: contact/causality exists, but the interaction or material/mechanism
  behaviour is physically impossible while it happens.
- 3: continuity, causality, and interaction are plausible, but the resulting
  state does not hold or change as visible forces allow.
- 4: no visible violation of continuity, causality, interaction, or persistence.
Allow ordinary occlusion, perspective/shadow changes, brief appearance flicker,
and transformations visibly explained by cutting, joining, consuming, adding,
pouring, spraying, erasing, or revealing. Do not penalize omitted task actions
as Physics failures. Do not invent an unseen process to explain a discontinuity.

For each dimension provide a numeric score (or Physics null), brief reasoning
based on visible observations, and up to three approximate original-source time
windows with evidence. Do not expose private chain-of-thought. Emit one
structured function call only.

Goal: <goal G>   [followed by the full video]
\end{lstlisting}
\paragraph{RBench first-person prompt (replaces the robot-specific wording).}
\begin{lstlisting}[style=promptstyle]
You are shown a 3x2 grid of six frames in chronological order (read row by
row). The frames come from a first-person video of a person interacting with objects.

Intended goal: <goal G>

Identify the hand and the goal-relevant objects from the goal and visible
frames. Evaluate the complete observed manipulation, without assuming that an
object is present or an action succeeded merely because the goal says so.

Score each aspect from 1 to 5. Be strict: 1 means failed, absent, or severely
implausible; 3 means broadly successful with visible problems; 5 means fully
successful and natural.

1. action_execution
   Did the visible hand perform the required action with coherent motion,
   contact, and ordering?
2. task_completion
   Does the observed final state satisfy the intended goal, including required
   intermediate state changes?
3. object_consistency
   Do the goal-relevant objects retain coherent identity, shape, count, and
   state across the frames?
4. hand_consistency
   Does the visible hand/body remain anatomically and temporally coherent,
   without duplication, disappearance, or impossible articulation?
5. physical_plausibility
   Are contact, support, gravity, containment, deformation, and object motion
   physically plausible, without penetration, floating, teleportation, or
   non-contact attachment?

Use the original RBench final-score convention: if action_execution or
task_completion is 1, total is 1; otherwise total is the mean of all five
aspect scores. Return strict JSON only, with this schema:
{
  "action_execution": {"reason": "...", "score": 1},
  "task_completion": {"reason": "...", "score": 1},
  "object_consistency": {"reason": "...", "score": 1},
  "hand_consistency": {"reason": "...", "score": 1},
  "physical_plausibility": {"reason": "...", "score": 1},
  "total": {"reason": "...", "score": 1.0}
}
\end{lstlisting}
